\documentclass[10pt]{article}
\usepackage[T1]{fontenc}
\usepackage[utf8]{inputenc}
\usepackage[french]{babel}
\usepackage{amsmath,amssymb,mathrsfs,mathpazo}
\usepackage{geometry}
\geometry{a4paper,margin=2cm}
\usepackage{graphicx,array,tabularx,tabularray,colortbl,multirow,multicol,booktabs,enumitem,pifont,ulem,fancybox,diagbox,cancel}
\usepackage{pstricks,pstricks-add,pst-plot,pst-tree,pst-node,pst-text,pst-func,pst-math,pst-bspline,pst-3dplot,pst-eucl,pst-solides3d}
\let\pscalculate\undefined
\usepackage{pst-bezier}
\usepackage{tikz,pgfplots,tkz-tab}
\usetikzlibrary{arrows,arrows.meta,patterns,calc,positioning,trees,decorations.markings,shapes.geometric,shapes.misc,shapes.arrows,intersections,decorations.pathmorphing,decorations.pathreplacing,backgrounds,shadows,babel}
\pgfplotsset{compat=1.18}
\usepgfplotslibrary{fillbetween}
\usepackage[np]{numprint}
\usepackage{listings,makecell,pas-tableur,twoopt}
\let\barre\undefined
\usepackage[pstricks]{tablvar}
\definecolor{darkgreen}{RGB}{0,100,0}
\providecommand{\red}{\color{red}}
\providecommand{\blue}{\color{blue}}
\providecommand{\green}{\color{green}}
\providecommand{\black}{\color{black}}
\providecommand{\white}{\color{white}}
\providecommand{\gray}{\color{gray}}
\providecommand{\cyan}{\color{cyan}}
\providecommand{\magenta}{\color{magenta}}
\providecommand{\yellow}{\color{yellow}}
\providecommand{\violet}{\color{violet}}
\providecommand{\orange}{\color{orange}}
\providecommand{\euro}{\texteuro}
\providecommand{\EUR}{\texteuro}
\providecommand{\ieme}{\textsuperscript{e}}
\providecommand{\degres}{\ensuremath{^\circ}}
\providecommand{\degree}{\ensuremath{^\circ}}
\providecommand{\vv}[1]{\overrightarrow{#1}}
\providecommand{\Coord}[1]{#1}
\providecommand{\decosix}{}
\providecommand{\starredbullet}{\ensuremath{\bullet}}
\pagestyle{empty}
\setlength{\parindent}{0pt}
\usepackage{adjustbox}\DeclareUnicodeCharacter{2236}{:}\DeclareUnicodeCharacter{2190}{\ensuremath{\leftarrow}}
\providecommand{\R}{\mathbb{R}}
\providecommand{\N}{\mathbb{N}}
\providecommand{\Z}{\mathbb{Z}}
\providecommand{\Q}{\mathbb{Q}}
\providecommand{\D}{\mathbb{D}}
\providecommand{\C}{\mathbb{C}}
\providecommand{\e}{{\mathrm{e}}}
\providecommand{\vect}[1]{\overrightarrow{#1}}
\providecommand{\vectt}[1]{\overrightarrow{\mathrm{#1}}}
\providecommand{\vvt}[1]{\overrightarrow{\mathrm{#1}}}
\providecommand{\np}[2][]{{#2~#1}}
\providecommand{\vv}[1]{{\overrightarrow{#1}}}
\renewcommand{\d}{{\mathrm{d}}}
\providecommand{\degres}{{^\circ}}
\providecommand{\degre}{{^\circ}}
\providecommand{\blue}{\color{blue}}
\providecommand{\red}{\color{red}}
\providecommand{\enskip}{\quad}
\providecommand{\llbracket}{\lbrack\!\lbrack}
\providecommand{\rrbracket}{\rbrack\!\rbrack}
\providecommand{\euro}{{\text{€}}}
\providecommand{\no}{{\text{n°}}}
\providecommand{\interval}[2]{{[#1;#2]}}
\providecommand{\intervFO}[2]{{[#1;#2[}}
\providecommand{\dotfill}{\ldots}

\providecommand{\ds}{\displaystyle}
\providecommand{\textding}[1]{#1}
\providecommand{\Oij}{\(\mathrm{O};\vec{\imath},\vec{\jmath}\)}
\providecommand{\Oijk}{\(\mathrm{O};\vec{\imath},\vec{\jmath},\vec{k}\)}
\providecommand{\Ouv}{\(\mathrm{O};\vec{u},\vec{v}\)}
\providecommand{\intervFF}[2]{\ensuremath{[#1;#2]}}
\renewcommand{\euro}{\ifmmode\text{\texteuro}\else\texteuro\fi}
\DeclareUnicodeCharacter{00B0}{\ifmmode^{\circ}\else\textdegree\fi}
\DeclareUnicodeCharacter{20AC}{\ifmmode\text{\texteuro}\else\texteuro\fi}
% Énoncés : Éducation nationale. Transcription et corrigés : APMEP, Denis Vergès et contributeurs. Découpage et classement : CoopMaths, Sébastien Lozano, Éric Elter et contributeurs. Adaptation de la présentation : Maths Explorer.
% CC BY-SA : https://coopmaths.fr/www/about/
% https://coopmaths.fr/alea/static/e3c/2020/tex/e3c_2020_00_sujet4_3.tex
\begin{document}
\noindent\textit{APMEP / CoopMaths. CC BY-SA. Présentation adaptée par Maths Explorer.}\par\medskip
\medskip

On dispose d'un paquet de cartes contenant un nombre identique de cartes de la catégorie \og Sciences\fg{} et de la catégorie \og Économie \fg. Une question liée à un de ces deux thèmes figure sur chaque carte.

Les cartes sont mélangées et on en tire une au hasard dans le paquet. Ensuite, on essaye de répondre à la question posée.

Un groupe de copains participe à ce jeu. Connaissant leurs points forts et leurs faiblesses, on estime qu'il a :

$\bullet~~$ 3 chances sur 4 de donner la bonne réponse lorsqu'il est interrogé en sciences;

$\bullet~~$ 1 chance sur 8 de donner la bonne réponse lorsqu'il est interrogé en économie.

\smallskip

On note $S$ l'évènement \og La question est dans la catégorie Sciences\fg{} et $B$ l'évènement \og La réponse donnée par le groupe est bonne \fg.

\bigskip

\textbf{Partie A :}

\medskip

\begin{enumerate}
\item Calculer $P(B \cap S)$.
\item Déterminer la probabilité que le groupe de copains réponde correctement à la
question posée.
\item Les évènements $S$ et $B$ sont-ils indépendants?
\end{enumerate}

\bigskip

\textbf{Partie B :}

\medskip

Pour participer à ce jeu, on doit payer 5~\euro{} de droit d'inscription. On recevra :

\setlength\parindent{1cm}
\begin{itemize}
\item[$\bullet$~~] 10 \euro{} si on est interrogé en sciences et que la réponse est correcte ;
\item[$\bullet$~~] 30 \euro{} si on est interrogé en économie et que la réponse est correcte ;
\item[$\bullet$~~] rien si la réponse donnée est fausse.
\end{itemize}
\setlength\parindent{0cm}

\smallskip

Soit $X$ la variable aléatoire qui, à chaque partie jouée, associe son gain. On appelle gain la différence en euros entre ce qui est reçu et les $5$~\euro{} de droit d'inscription.

\medskip

\begin{enumerate}
\item Déterminer la loi de probabilité de $X$.
\item Que retourne la fonction Jeu écrite ci-dessous en langage Python
avec les listes : 

L =[-5~;~5~;~25] et G = [0,5625~;~0,375~;~\np{0,0625}] ?

\begin{center}

\begin{tabularx}{0.4\linewidth}{|X|}\hline
def \texttt{Jeu(L,G):}\\
\quad  \texttt{n = len(L)}\\
\quad \texttt{E = 0}\\
\quad \texttt{for i in range(n):}\\
\qquad \texttt{E =E + L[i]*G[i]} \\
\quad \texttt{return(E)}\\ \hline
\end{tabularx}
\end{center}
\end{enumerate}

\bigskip

\end{document}
